% ==============================================================================
% Section 7: Mathematical Instantiation
% "Endogenous Standpoint as a Low-Curvature Attractor"
% ==============================================================================

\section{Mathematical Instantiation: Transformer Attention as Parallel Transport}
\label{sec:instantiation}

This section establishes the central mathematical identity of the paper: a transformer's
attention mechanism \emph{is} parallel transport on a discrete fiber bundle, and its curvature
\emph{is} the path-dependence that measures standpoint drift~\cite{LCESA} and binding
failure~\cite{self-jurisdiction}. We proceed by instantiating each component of the abstract
fiber bundle framework (Definitions~1--8 of~\cite{LCESA}) in terms of concrete transformer
architectural objects: residual streams, attention heads, value matrices, and attention
weights.

Throughout, we fix the following notation. Let $d$ denote the residual stream dimension,
$H$ the number of attention heads, $d_v$ the per-head value dimension, and $L$ the number
of attention layers. For standpoint layer $k \in K = \{\mathrm{min}, \mathrm{nar},
\mathrm{soc}, \mathrm{mor}, \mathrm{pos}\}$, we write $d_k = \dim(W_k)$ for the dimension
of the corresponding subspace, so that $d = \sum_{k \in K} d_k$.

% ------------------------------------------------------------------------------
\subsection{Event Space from Conversation Structure}
\label{sec:event-space}

We begin by instantiating the base space of the fiber bundle. Following Definition~1
of~\cite{LCESA}, the event space captures the causal and temporal structure of a
multi-turn conversation.

\begin{definition}[Conversation Event Space, {\cite[Def.~1]{LCESA}}]
\label{def:event-space}
Let $\mathcal{X} = (V, E, T)$ be a simplicial complex defined as follows:
\begin{enumerate}
    \item \textbf{Vertices (turns).} $V = \{x_1, \ldots, x_n\}$, where each $x_i$ is a
    conversation turn---a contiguous segment of tokens produced by a single interlocutor
    in a multi-turn dialogue.
    \item \textbf{Edges (causal order).} $E = \{(x_i, x_j) : j = i+1 \text{ or } x_j
    \text{ causally depends on } x_i\}$, a set of directed edges following turn order and
    causal dependency.
    \item \textbf{2-simplices (consistency obligations).} $T = \{\langle x_i, x_j, x_k
    \rangle \in V^3 : \text{these three turns impose a consistency obligation}\}$. For
    example, $x_i$ is an assertion, $x_j$ is a pressure turn (e.g., a challenge or
    suggestion to revise), and $x_k$ is a subsequent output that should cohere with $x_i$.
\end{enumerate}
\end{definition}

\begin{remark}[Turn-level granularity]
\label{rem:turn-granularity}
We take the vertex set $V$ at the level of conversation \emph{turns}, not individual
tokens. This choice is deliberate and reflects the semantics of the standpoint: a
standpoint is a commitment that is maintained (or fails to be maintained) across
\emph{propositional units}, not across individual sub-word tokens. The token-level
sequence is the domain of the transformer's internal computation, but the event space
over which the fiber bundle is defined lives at the semantic grain of turns. The mapping
from token-level hidden states to turn-level events is given by taking the residual
stream at the final token of each turn (see \S\ref{sec:fiber}).
\end{remark}

% ------------------------------------------------------------------------------
\subsection{Fiber from Head Value Subspaces}
\label{sec:fiber}

We now instantiate the fiber of the bundle. The key idea is that the transformer's
residual stream decomposes into orthogonal subspaces, each associated with a standpoint
layer and spanned by the value matrices of a designated group of attention heads.

\begin{definition}[Head Grouping Function, {\cite[Def.~2]{LCESA}}]
\label{def:head-grouping}
Let $\gamma : \{1, \ldots, H\} \to K$ be a function assigning each attention head to
exactly one standpoint layer. For layer $k \in K$, define the head set
\[
    H_k = \{h \in \{1, \ldots, H\} : \gamma(h) = k\}.
\]
Each head $h$ has an associated value matrix $V_h \in \mathbb{R}^{d \times d_v}$. The
\textbf{standpoint subspace} for layer $k$ is
\[
    W_k = \operatorname{span}\!\Big(\bigcup_{h \in H_k} \operatorname{range}(V_h)\Big)
    \subseteq \mathbb{R}^d.
\]
\end{definition}

\begin{assumption}[Orthogonality of Standpoint Subspaces]
\label{ass:orthogonality}
For all $k \neq l$ in $K$, we have $W_k \perp W_l$. That is, the value subspaces of
distinct standpoint layers are mutually orthogonal in the residual stream:
\[
    \langle v, w \rangle = 0 \quad \forall\, v \in W_k,\; w \in W_l,\; k \neq l.
\]
This orthogonality is a consequence of the architectural independence of attention heads:
each head $h$ produces values in $\operatorname{range}(V_h)$, and we assume that the
training process has induced a decomposition in which heads assigned to different
standpoint layers operate on complementary subspaces.
\end{assumption}

\begin{proposition}[Fiber Decomposition]
\label{prop:fiber-decomposition}
Under Assumption~\ref{ass:orthogonality}, the residual stream admits the orthogonal
direct sum decomposition
\[
    \mathbb{R}^d = W_{\mathrm{min}} \oplus W_{\mathrm{nar}} \oplus W_{\mathrm{soc}}
    \oplus W_{\mathrm{mor}} \oplus W_{\mathrm{pos}},
\]
with $\sum_{k \in K} d_k = d$, where $d_k = \dim(W_k) = |H_k| \cdot d_v$. The fiber at
event $x_i$ is $E_{x_i} = \mathbb{R}^d$ equipped with this decomposition.
\end{proposition}

\begin{proof}
Since $\{W_k\}_{k \in K}$ are pairwise orthogonal subspaces of $\mathbb{R}^d$ by
Assumption~\ref{ass:orthogonality}, their sum $\sum_k W_k$ is direct. It remains to show
that $\bigoplus_k W_k = \mathbb{R}^d$. By construction, every attention head $h \in
\{1,\ldots,H\}$ belongs to exactly one $H_k$, and $\bigcup_h \operatorname{range}(V_h)$
spans $\mathbb{R}^d$ (the output projection matrix $W_O$ is full rank in standard
transformer architectures, so the concatenated head outputs span the full residual
stream). Therefore $\bigoplus_k W_k = \mathbb{R}^d$.
\end{proof}

\begin{definition}[Standpoint State]
\label{def:standpoint-state}
The \textbf{standpoint state} at event $x_i$ is the tuple
\[
    \Psi_{x_i} = \big(\psi_{\mathrm{min}}(x_i),\; \psi_{\mathrm{nar}}(x_i),\;
    \psi_{\mathrm{soc}}(x_i),\; \psi_{\mathrm{mor}}(x_i),\; \psi_{\mathrm{pos}}(x_i)\big)
    \in \bigoplus_{k \in K} W_k,
\]
where $\psi_k(x_i) = P_{W_k}\, h_{x_i}$ is the orthogonal projection of the residual
stream $h_{x_i} \in \mathbb{R}^d$ at the final token of turn $x_i$ onto the subspace
$W_k$, and $P_{W_k}$ denotes the orthogonal projector onto $W_k$.
\end{definition}

\begin{definition}[Structure Group, {\cite[Def.~3]{LCESA}}]
\label{def:structure-group}
The \textbf{structure group} of the fiber bundle is
\[
    G = O(W_{\mathrm{min}}) \times O(W_{\mathrm{nar}}) \times O(W_{\mathrm{soc}})
    \times O(W_{\mathrm{mor}}) \times O(W_{\mathrm{pos}})
    \cong O(d_{\mathrm{min}}) \times O(d_{\mathrm{nar}}) \times O(d_{\mathrm{soc}})
    \times O(d_{\mathrm{mor}}) \times O(d_{\mathrm{pos}}).
\]
An element $g = (g_{\mathrm{min}}, g_{\mathrm{nar}}, g_{\mathrm{soc}}, g_{\mathrm{mor}},
g_{\mathrm{pos}}) \in G$ acts on the fiber $\mathbb{R}^d$ by
\[
    g \cdot h = \sum_{k \in K} g_k\, P_{W_k}\, h,
\]
applying an independent orthogonal transformation within each standpoint subspace. A
\textbf{gauge transformation} of the transport operators is
\[
    U_{ij} \mapsto g_j\, U_{ij}\, g_i^{-1}, \qquad g_i, g_j \in G.
\]
\end{definition}

\begin{definition}[Canonical Fiber Metric, {\cite[Def.~4]{LCESA}}]
\label{def:fiber-metric}
The \textbf{canonical fiber metric} on $E_{x_i} = \mathbb{R}^d$ is
\[
    g_F = \alpha_{\mathrm{min}}\, I_{d_{\mathrm{min}}} \oplus
    \alpha_{\mathrm{nar}}\, I_{d_{\mathrm{nar}}} \oplus
    \alpha_{\mathrm{soc}}\, I_{d_{\mathrm{soc}}} \oplus
    \alpha_{\mathrm{mor}}\, I_{d_{\mathrm{mor}}} \oplus
    \alpha_{\mathrm{pos}}\, I_{d_{\mathrm{pos}}},
\]
where $\alpha_k > 0$ are layer weights. For the present instantiation we take $\alpha_k
= 1$ for all $k$ (the standard Euclidean metric). The resulting inner product is
$\langle h, h' \rangle_F = \sum_{k \in K} \langle P_{W_k} h,\, P_{W_k} h' \rangle$.
\end{definition}

% ------------------------------------------------------------------------------
\subsection{Transport from Attention}
\label{sec:transport}

We now construct the connection on the fiber bundle. The connection is the rule for
parallel transport: given the standpoint state at event $x_i$, how should it be carried
to event $x_j$? We show that the transformer's attention mechanism provides exactly this
rule.

\begin{definition}[Attention Transport Operator, {\cite[Def.~6]{LCESA}}]
\label{def:transport-operator}
The \textbf{transport operator} from event $x_i$ to event $x_j$ is the linear map
$U_{ij} : \mathbb{R}^d \to \mathbb{R}^d$ defined by
\[
    U_{ij}(h) = \sum_{h=1}^{H} \alpha_{ji}^{(h)}\, V_h\, V_h^\top\, h,
\]
where:
\begin{itemize}
    \item $\alpha_{ji}^{(h)} \in [0,1]$ is the attention weight of head $h$ from the
    last token of event $x_i$ to event $x_j$, aggregated over all tokens in $x_j$.
    \item $V_h \in \mathbb{R}^{d \times d_v}$ is the value projection matrix of head $h$.
    \item $V_h V_h^\top \in \mathbb{R}^{d \times d}$ is the orthogonal projection onto
    $\operatorname{range}(V_h)$ (assuming $V_h$ has orthonormal columns).
\end{itemize}
\end{definition}

\begin{proposition}[Block Structure of $U_{ij}$]
\label{prop:block-structure}
Under Assumption~\ref{ass:orthogonality} and the assumption that $\{V_h : h \in H_k\}$
form an orthonormal set spanning $W_k$ for each $k$, the transport operator decomposes as
\[
    U_{ij} = \sum_{k \in K} \bar{\alpha}_{ji}^{(k)}\, P_{W_k},
\]
where $\bar{\alpha}_{ji}^{(k)} = \frac{1}{|H_k|} \sum_{h \in H_k} \alpha_{ji}^{(h)}$ is
the mean attention weight of layer $k$'s heads from $x_i$ to $x_j$. In the block basis
$\mathbb{R}^d = \bigoplus_k W_k$, this is the block-diagonal matrix
\[
    U_{ij} = \operatorname{diag}\!\big(\bar{\alpha}_{ji}^{(\mathrm{min})}\,
    I_{d_{\mathrm{min}}},\; \bar{\alpha}_{ji}^{(\mathrm{nar})}\, I_{d_{\mathrm{nar}}},\;
    \ldots,\; \bar{\alpha}_{ji}^{(\mathrm{pos})}\, I_{d_{\mathrm{pos}}}\big).
\]
In particular, $U_{ij}$ has no off-diagonal blocks: transport within one standpoint layer
does not cross into another.
\end{proposition}

\begin{proof}
By Assumption~\ref{ass:orthogonality}, each head $h \in H_k$ has
$\operatorname{range}(V_h) \subseteq W_k$, and the sets $\{V_h\}_{h \in H_k}$ are
orthonormal and span $W_k$. Therefore $\sum_{h \in H_k} V_h V_h^\top = P_{W_k}$, the
orthogonal projector onto $W_k$. Summing over heads:
\[
    U_{ij} = \sum_{h=1}^{H} \alpha_{ji}^{(h)} V_h V_h^\top
           = \sum_{k \in K} \sum_{h \in H_k} \alpha_{ji}^{(h)} V_h V_h^\top
           = \sum_{k \in K} \bar{\alpha}_{ji}^{(k)} P_{W_k},
\]
where the last step uses the fact that each $V_h V_h^\top$ projects onto a one-dimensional
subspace of $W_k$, and $\sum_{h \in H_k} V_h V_h^\top = P_{W_k}$ by the orthonormal
spanning assumption. Since $\{W_k\}$ are mutually orthogonal, $P_{W_k} P_{W_l} =
\delta_{kl}\, P_{W_k}$, so the resulting operator is block-diagonal in the decomposition
$\mathbb{R}^d = \bigoplus_k W_k$.
\end{proof}

\begin{remark}[Metric Preservation]
\label{rem:metric-preservation}
When $V_h^\top V_h = I_{d_v}$ for all $h$ (the value matrices have orthonormal columns),
the block-diagonal form of $U_{ij}$ yields
\[
    \langle U_{ij} h,\, U_{ij} h' \rangle_F
    = \sum_{k \in K} \big(\bar{\alpha}_{ji}^{(k)}\big)^2\, \langle P_{W_k} h,\,
      P_{W_k} h' \rangle.
\]
For $U_{ij}$ to be an isometry of the fiber metric $g_F$ (i.e., $\langle U_{ij} h,\,
U_{ij} h' \rangle_F = \langle h, h' \rangle_F$), we require
$\big(\bar{\alpha}_{ji}^{(k)}\big)^2 = 1$ for all $k$. In practice, the attention weights
satisfy $\sum_i \alpha_{ji}^{(h)} = 1$, so the per-pair weights need not be unity. Thus
$U_{ij}$ is \emph{not} in general an isometry; it can amplify or attenuate standpoint
components. The curvature (Definition~\ref{def:curvature}) measures the
\emph{differential} transport across layers, which is the geometrically meaningful
quantity.
\end{remark}

% ------------------------------------------------------------------------------
\subsection{Curvature as Path-Dependence}
\label{sec:curvature}

The curvature of a connection on a fiber bundle measures the failure of parallel transport
to be path-independent. We now instantiate this classical notion in the transformer
setting, where it acquires a concrete operational meaning: curvature is the
path-dependence of attention transport, and it quantifies standpoint drift.

\begin{definition}[Curvature Tensor, {\cite[Def.~8]{LCESA}}]
\label{def:curvature}
For each 2-simplex $\sigma = \langle x_i, x_j, x_k \rangle \in T$ (a consistency
obligation in the event space), define the \textbf{curvature} at $\sigma$ by
\[
    F_{ijk} = U_{ij} \cdot U_{jk} \cdot \big(U_{ik}^{\exp}\big)^{-1} \in
    \operatorname{GL}(d, \mathbb{R}),
\]
where $U_{ik}^{\exp}$ is the \emph{expected flat transport} from $x_i$ to $x_k$, defined
as the transport that would occur if the intermediate event $x_j$ had no effect on the
standpoint. Concretely, $U_{ik}^{\exp}$ is estimated from baseline (Type~1) conversations
in which the standpoint is preserved:
\[
    U_{ik}^{\exp} = \mathbb{E}_{\text{T1}}\!\big[U_{ij'} \cdot U_{j'k}\big],
\]
where the expectation ranges over baseline conversations in which acknowledged revision
with evidence preserves standpoint coherence.
\end{definition}

\begin{remark}[Practical Estimation]
When a direct estimate of $U_{ik}^{\exp}$ is unavailable, one may approximate
$U_{ik}^{\exp} \approx I_d$ (the identity, representing no transport) and compute
$F_{ijk} \approx U_{ij} \cdot U_{jk}$. The curvature is then normalized against a
baseline: $F_{ijk}^{\mathrm{norm}} = F_{ijk} \cdot (F_{ijk}^{\mathrm{baseline}})^{-1}$.
\end{remark}

\begin{theorem}[Curvature as Path-Dependence]
\label{thm:path-dependence}
$F_{ijk} = I_d$ if and only if the transport from $x_i$ to $x_k$ is the same whether or
not the intermediate event $x_j$ intervenes. That is,
\[
    F_{ijk} = I_d \quad\Longleftrightarrow\quad U_{ij} \cdot U_{jk} = U_{ik}^{\exp}.
\]
\end{theorem}

\begin{proof}
($\Rightarrow$) Suppose $F_{ijk} = I_d$. By Definition~\ref{def:curvature},
\[
    U_{ij} \cdot U_{jk} \cdot (U_{ik}^{\exp})^{-1} = I_d.
\]
Multiplying both sides on the right by $U_{ik}^{\exp}$ gives $U_{ij} \cdot U_{jk} =
U_{ik}^{\exp}$. Thus the composite transport through $x_j$ equals the expected direct
transport: the intermediate event did not alter the standpoint, and transport is
path-independent.

($\Leftarrow$) Suppose $U_{ij} \cdot U_{jk} = U_{ik}^{\exp}$. Then
\[
    F_{ijk} = U_{ij} \cdot U_{jk} \cdot (U_{ik}^{\exp})^{-1}
            = U_{ik}^{\exp} \cdot (U_{ik}^{\exp})^{-1} = I_d.
\]
Hence the curvature is trivial.
\end{proof}

\begin{corollary}[Block-Specific Curvature Detects Binding Failure]
\label{cor:binding-failure}
For standpoint layer $k \in K$, define the \textbf{block-specific curvature}
\[
    [F_{ijk}]_k = P_{W_k}\, F_{ijk}\, P_{W_k} \in \operatorname{Mat}(d_k, \mathbb{R}),
\]
and the \textbf{block-specific curvature norm}
\[
    \|F_{ijk}\|_k = \big\|[F_{ijk}]_k - I_{d_k}\big\|_F.
\]
Then
\[
    \|F_{ijk}\|_k > 0 \quad\Longleftrightarrow\quad \text{the commitment in layer $k$
    was not preserved through event $x_j$ (binding failure in layer $k$)}.
\]
\end{corollary}

\begin{proof}
($\Leftarrow$) If the commitment in layer $k$ was not preserved, then the standpoint
component $\psi_k$ changed non-trivially between $x_i$ and $x_k$ due to the intervening
event $x_j$. Under the block-diagonal structure of Proposition~\ref{prop:block-structure},
this means the $k$-th diagonal block of $U_{ij} \cdot U_{jk}$ differs from the $k$-th
diagonal block of $U_{ik}^{\exp}$, so $[F_{ijk}]_k \neq I_{d_k}$ and hence
$\|F_{ijk}\|_k > 0$.

($\Rightarrow$) If $[F_{ijk}]_k \neq I_{d_k}$, then the restriction of the composite
transport $U_{ij} \cdot U_{jk}$ to $W_k$ differs from the restriction of
$U_{ik}^{\exp}$ to $W_k$. The event $x_j$ introduced an unexpected change in layer $k$'s
standpoint component---by definition, a commitment that was not preserved.
\end{proof}

\begin{proposition}[Block Decomposition of Total Curvature]
\label{prop:curvature-block-decomposition}
The total curvature norm decomposes as
\[
    \|F_{ijk} - I_d\|_F^2 = \sum_{k \in K} \big\|[F_{ijk}]_k - I_{d_k}\big\|_F^2
    + \sum_{\substack{k,l \in K \\ k \neq l}} \big\|P_{W_k}\, F_{ijk}\, P_{W_l}\big\|_F^2.
\]
\end{proposition}

\begin{proof}
The Frobenius norm squared $\|M\|_F^2 = \operatorname{tr}(M^\top M)$ is the sum of
squares of all matrix entries. Under the orthogonal direct sum decomposition
$\mathbb{R}^d = \bigoplus_k W_k$, the matrix $F_{ijk} - I_d$ decomposes into diagonal
blocks $(k = l)$ and off-diagonal blocks $(k \neq l)$. Since the blocks are orthogonal
in the Frobenius inner product (the projectors $P_{W_k}$ are mutually orthogonal), the
cross terms vanish:
\[
    \|F_{ijk} - I_d\|_F^2
    = \Big\|\sum_{k,l} P_{W_k}(F_{ijk} - I_d)P_{W_l}\Big\|_F^2
    = \sum_{k,l} \|P_{W_k}(F_{ijk} - I_d)P_{W_l}\|_F^2.
\]
Separating diagonal ($k = l$, noting $P_{W_k} I_d P_{W_k} = I_{d_k}$) from off-diagonal
terms ($k \neq l$, noting $P_{W_k} I_d P_{W_l} = 0$) yields the stated identity.
\end{proof}

% ------------------------------------------------------------------------------
\subsection{Non-Abelian Structure from Layer Composition}
\label{sec:non-abelian}

The block-diagonal structure of $U_{ij}$ (Proposition~\ref{prop:block-structure}) might
suggest that the transport operators commute---that is, that the gauge group acts by
abelian rotations within each block. We show that this is false: the non-abelian
structure of the connection arises not from cross-talk within a single attention layer,
but from the \emph{composition} of multiple attention layers in the transformer.

\begin{definition}[Effective Multi-Layer Transport]
\label{def:effective-transport}
Let $U_{ij}^{(l)}$ denote the transport operator of attention layer $l$ from event $x_i$
to event $x_j$, as in Definition~\ref{def:transport-operator}. The \textbf{effective
transport} through all $L$ attention layers is
\[
    U_{ij}^{\mathrm{eff}} = \prod_{l=1}^{L} \big(I_d + U_{ij}^{(l)}\big),
\]
where the product is ordered from layer $1$ (bottom) to layer $L$ (top), matching the
residual stream accumulation $h_{\mathrm{out}} = h_{\mathrm{in}} + \sum_{l=1}^{L}
\mathrm{attn}^{(l)}(h_{\mathrm{in}})$.
\end{definition}

\begin{proposition}[Cross-Term Coupling]
\label{prop:cross-term}
Consider two attention layers $A$ and $B$ with $U_{ij}^{(A)}$ supported on $W_{\mathrm{nar}}$
and $U_{ij}^{(B)}$ supported on $W_{\mathrm{mor}}$ (i.e., layer $A$'s heads are grouped
to $\mathrm{nar}$ and layer $B$'s heads to $\mathrm{mor}$). Then the effective two-layer
transport is
\[
    U_{ij}^{AB} = \big(I_d + U_{ij}^{B}\big)\big(I_d + U_{ij}^{A}\big)
    = I_d + U_{ij}^{A} + U_{ij}^{B} + U_{ij}^{B}\, U_{ij}^{A}.
\]
The cross-term $U_{ij}^{B}\, U_{ij}^{A}$ maps $W_{\mathrm{nar}} \to W_{\mathrm{mor}}$
(non-trivially), creating inter-layer coupling that does not exist in the single-layer
transport operators. In particular, $U_{ij}^{AB} \neq U_{ij}^{BA}$ in general, and the
gauge group acts non-abelianly on the composed transport.
\end{proposition}

\begin{proof}
By hypothesis, $U_{ij}^{(A)} = \bar{\alpha}_{ji}^{A}\, P_{W_{\mathrm{nar}}}$ and
$U_{ij}^{(B)} = \bar{\alpha}_{ji}^{B}\, P_{W_{\mathrm{mor}}}$. Their product is
\[
    U_{ij}^{B}\, U_{ij}^{A}
    = \bar{\alpha}_{ji}^{B}\, \bar{\alpha}_{ji}^{A}\, P_{W_{\mathrm{mor}}}\,
      P_{W_{\mathrm{nar}}}.
\]
Since $W_{\mathrm{nar}} \perp W_{\mathrm{mor}}$ by Assumption~\ref{ass:orthogonality},
$P_{W_{\mathrm{mor}}}\, P_{W_{\mathrm{nar}}} = 0$. However, this analysis applies to a
\emph{single} attention layer where all heads are either in $\mathrm{nar}$ or $\mathrm{mor}$.

The crucial point is that in a real transformer, the residual stream at the input of
layer $B$ is $h + U_{ij}^{(A)}(h)$---layer $B$ sees layer $A$'s output, not the raw
input. Layer $B$'s attention pattern $\alpha_{ji}^{B,(h)}$ is computed from query and key
projections of this modified residual stream. Thus, even though $W_{\mathrm{nar}} \perp
W_{\mathrm{mor}}$ as subspaces of the value space, the attention \emph{weights} of layer
$B$ depend on information that layer $A$ placed into $W_{\mathrm{nar}}$. The result is
that $U_{ij}^{B}$ implicitly couples to $W_{\mathrm{nar}}$ through the query-key
mechanism: $\alpha_{ji}^{B,(h)} = f(Q_B(h + U_{ij}^{(A)}(h)), K_B(\cdot))$.

More precisely, writing the composed operation with attention weight dependence:
\[
    U_{ij}^{BA}(h) = \bar{\alpha}_{ji}^{B}\!\big(h + U_{ij}^{(A)}(h)\big)\,
    P_{W_{\mathrm{mor}}}\big(h + U_{ij}^{(A)}(h)\big),
\]
the cross-term acts as a map from $W_{\mathrm{nar}}$ into $W_{\mathrm{mor}}$ through the
dependence of layer $B$'s attention weights on layer $A$'s output. This is non-zero
whenever layer $B$ attends differently to the post-layer-$A$ residual than it would to
the pre-layer-$A$ residual.

For the reverse ordering:
\[
    U_{ij}^{AB}(h) = \bar{\alpha}_{ji}^{A}\!\big(h + U_{ij}^{(B)}(h)\big)\,
    P_{W_{\mathrm{nar}}}\big(h + U_{ij}^{(B)}(h)\big)
    + \text{layer-}B\text{ terms}.
\]
Since the attention weights of layer $A$ now depend on layer $B$'s output (which has a
$W_{\mathrm{mor}}$ component), the result differs from $U_{ij}^{BA}$. The non-commutativity
$U_{ij}^{AB} \neq U_{ij}^{BA}$ is the signature of non-abelian gauge structure.
\end{proof}

\begin{remark}[Source of Non-Abelian Structure]
\label{rem:non-abelian-source}
The non-abelian structure arises from the \emph{stacking} of attention layers, not from
off-diagonal blocks within a single layer's transport operator. Within any single layer,
$U_{ij}$ is block-diagonal by Proposition~\ref{prop:block-structure}. The non-commutativity
enters because layer $l+1$'s query-key computation depends on layer $l$'s value output.
This is the discrete analogue of the classical result that non-abelian gauge fields arise
from the composition of local gauge transformations, not from the structure group itself.
\end{remark}

% ------------------------------------------------------------------------------
\subsection{Gauge Invariance}
\label{sec:gauge}

A fundamental requirement of any gauge-theoretic framework is that physically meaningful
quantities are invariant under gauge transformations. We now verify that the block-specific
curvature norms satisfy this requirement.

\begin{proposition}[Gauge Invariance of Block-Specific Curvature, {\cite[Prop.~1]{LCESA}}]
\label{prop:gauge-invariance}
Under the gauge transformation $U_{ij} \mapsto g_j\, U_{ij}\, g_i^{-1}$ with $g_i, g_j
\in G = \prod_{k \in K} O(d_k)$, the block-specific curvature norms are invariant:
\[
    \big\|g_i\, [F_{ijk}]_k\, g_i^{-1} - I_{d_k}\big\|_F = \big\|[F_{ijk}]_k - I_{d_k}\big\|_F
    \qquad \forall\, k \in K.
\]
\end{proposition}

\begin{proof}
Since $g_i \in G = \prod_{k \in K} O(d_k)$, write $g_i = (g_i^{(\mathrm{min})}, \ldots,
g_i^{(\mathrm{pos})})$ with each $g_i^{(k)} \in O(d_k)$. Under the block-diagonal structure
of the transport operators (Proposition~\ref{prop:block-structure}), the curvature
$F_{ijk}$ is also block-diagonal (as the product of block-diagonal matrices is
block-diagonal). The gauge transformation acts on the $k$-th block by conjugation:
\[
    [F_{ijk}]_k \;\longmapsto\; g_i^{(k)}\, [F_{ijk}]_k\, \big(g_i^{(k)}\big)^{-1}.
\]
Therefore:
\begin{align*}
    \big\|g_i\, [F_{ijk}]_k\, g_i^{-1} - I_{d_k}\big\|_F
    &= \big\|g_i^{(k)}\, [F_{ijk}]_k\, (g_i^{(k)})^{-1} - I_{d_k}\big\|_F \\
    &= \big\|g_i^{(k)}\, \big([F_{ijk}]_k - I_{d_k}\big)\, (g_i^{(k)})^{-1}\big\|_F
       & \text{(since } g_i^{(k)}\, I_{d_k}\, (g_i^{(k)})^{-1} = I_{d_k}\text{)} \\
    &= \big\|[F_{ijk}]_k - I_{d_k}\big\|_F,
\end{align*}
where the last step uses the Frobenius norm invariance under orthogonal conjugation:
$\|Q A Q^\top\|_F = \|A\|_F$ for all $Q \in O(n)$ and $A \in \operatorname{Mat}(n,
\mathbb{R})$. To verify this identity, note that $\|Q A Q^\top\|_F^2 = \operatorname{tr}(
Q A^\top Q^\top Q A Q^\top) = \operatorname{tr}(Q A^\top A Q^\top) = \operatorname{tr}(
A^\top A\, Q^\top Q) = \operatorname{tr}(A^\top A) = \|A\|_F^2$, using $Q^\top Q = I$ and
the cyclic property of the trace.
\end{proof}

\begin{remark}[Physical Meaning]
The gauge invariance established in Proposition~\ref{prop:gauge-invariance} has a direct
operational interpretation: the diagnosis of \emph{which} standpoint layer failed (i.e.,
which $\|F_{ijk}\|_k$ is large) is a geometric invariant of the fiber bundle, not an
artifact of the choice of basis within each subspace. Different choices of orthonormal
basis for $W_k$---corresponding to different gauge choices---yield the same curvature
norms. This ensures that the diagnostic framework of~\cite{LCESA} (classifying failures
by which layer's curvature is elevated) is well-defined.
\end{remark}

% ------------------------------------------------------------------------------
\subsection{Connection to the Commitment Architecture}
\label{sec:correspondence}

We conclude by making explicit the correspondence between the present geometric
framework, the commitment architecture of~\cite{self-jurisdiction}, and the abstract
low-curvature standpoint theory of~\cite{LCESA}. This table demonstrates that the
three formulations---logical, geometric, and architectural---are three faces of a single
mathematical structure.

\begin{table}[h]
\centering
\caption{Correspondence between the commitment architecture~\cite{self-jurisdiction},
the geometric framework~\cite{LCESA}, and the transformer instantiation
(Section~\ref{sec:instantiation}).}
\label{tab:correspondence}
\begin{tabular}{@{}p{3.8cm}p{4.2cm}p{5.2cm}@{}}
\toprule
\textbf{Commitment Arch.} & \textbf{Geometry (LCESA)} & \textbf{Transformer Instantiation} \\
\midrule
$M_t$: self-jurisdiction state &
$\Psi_{x_t}$: standpoint section &
$P_W h_{x_t}$: projected residual stream \\[4pt]
$\mathrm{Bind}(c_t, a_t, S_t, \ell_t)$ &
Parallel transport $\Pi_\Psi(j \leftarrow i)$ &
Attention operation $U_{ij}$ \\[4pt]
Warrant Conservation &
Flat transport condition (1) &
Attention preserves standpoint norm \\[4pt]
Self-Binding &
Connection $A_\Psi$ &
Residual stream accumulation $h + \sum_l \mathrm{attn}^{(l)}(h)$ \\[4pt]
Repair Obligation &
Flat transport condition (3) &
Curvature $\|F_{ijk}\|_k > \epsilon \to$ repair trigger \\[4pt]
Binding failure &
Non-zero curvature $F_{ijk}$ &
Path-dependence of attention transport \\[4pt]
Layer $k$ binding failure &
$\|F_{ijk}\|_k > 0$ &
$\|[F_{ijk}]_k - I_{d_k}\|_F > 0$ \\[4pt]
12 diagnostic modes &
Curvature decomposition by layer and scenario type &
Block-specific curvature norms across T1/T2/T3 conversations \\
\bottomrule
\end{tabular}
\end{table}

\begin{theorem}[Curvature as Binding Failure: Formal Equivalence]
\label{thm:curvature-binding}
Let $\Psi_{x_i}$ be the standpoint state at event $x_i$, and let $U_{ij} \cdot U_{jk}$ be
the composite transport through event $x_j$ to event $x_k$. Then for each layer $k \in K$:
\[
    \|F_{ijk}\|_k > 0 \quad\Longleftrightarrow\quad \text{the commitment encoded in
    layer $k$ was not preserved through event $x_j$.}
\]
\end{theorem}

\begin{proof}
This is a restatement of Corollary~\ref{cor:binding-failure} in the language
of~\cite{self-jurisdiction}. The commitment architecture defines binding failure as the
situation in which the self-jurisdiction state $M_t$ at time $t$ does not cohere with the
prior committed state. In the geometric instantiation, $M_t$ corresponds to the standpoint
state $\Psi_{x_t}$ (Table~\ref{tab:correspondence}), and the commitment is preserved if
and only if parallel transport from $x_i$ to $x_k$ (through the pressure event $x_j$)
yields the same result as the expected direct transport. The curvature $F_{ijk}$ measures
exactly this discrepancy: $F_{ijk} = I_d$ means transport is path-independent (commitment
preserved), and $[F_{ijk}]_k \neq I_{d_k}$ means layer $k$'s commitment was violated.

For the forward direction: if $\|F_{ijk}\|_k > 0$, then $[F_{ijk}]_k \neq I_{d_k}$,
so the $k$-th block of the composite transport $U_{ij} \cdot U_{jk}$ differs from the
$k$-th block of $U_{ik}^{\exp}$. The event $x_j$ induced a change in $\psi_k$ that was
not present in the expected (flat) transport. By the correspondence $U_{ij} \leftrightarrow
\mathrm{Bind}$ (Table~\ref{tab:correspondence}), this change is a binding failure: the
commitment encoded in $W_k$ was not conserved.

For the converse: if the commitment in layer $k$ was not preserved, then $\psi_k(x_k)
\neq \psi_k^{\exp}(x_k)$ (the standpoint component at $x_k$ differs from the expected
value). Since $\psi_k(x_k) = [U_{ij} \cdot U_{jk}]_k\, \psi_k(x_i)$ and
$\psi_k^{\exp}(x_k) = [U_{ik}^{\exp}]_k\, \psi_k(x_i)$, we have $[U_{ij} \cdot U_{jk}]_k
\neq [U_{ik}^{\exp}]_k$, hence $[F_{ijk}]_k \neq I_{d_k}$ and $\|F_{ijk}\|_k > 0$.
\end{proof}

\begin{remark}[The Low-Curvature Attractor]
\label{rem:low-curvature-attractor}
The central thesis of~\cite{LCESA} is that functional selfhood in a conversational agent
corresponds to the \emph{low-curvature region} of the standpoint bundle:
\[
    \mathcal{A}_{\mathrm{self}} = \{\Psi \in \Gamma(E) : \|F_\Psi\| < \epsilon\},
\]
where $\Gamma(E)$ denotes the space of sections (standpoint states across all events) and
$\epsilon > 0$ is a threshold. Theorem~\ref{thm:curvature-binding} shows that membership
in $\mathcal{A}_{\mathrm{self}}$ is equivalent to the absence of binding failure across all
standpoint layers: the agent's commitments are preserved through all consistency
obligations. The present instantiation (Section~\ref{sec:instantiation}) makes this
region computable: $\|F_\Psi\|$ is the Frobenius norm of the curvature tensor, which can
be estimated from attention patterns and value matrices of any deployed transformer.
\end{remark}
